\documentclass[11pt]{article}
%DIF LATEXDIFF DIFFERENCE FILE
%DIF DEL main_tables_figures_before.tex   Sun Sep 20 19:58:41 2026
%DIF ADD main_tables_figures_after.tex    Sun Sep 20 19:58:42 2026
\usepackage[margin=1in]{geometry}
\usepackage[T1]{fontenc}
\usepackage[utf8]{inputenc}
\usepackage{mathptmx,amsmath,amssymb,amsthm,mathtools}
\usepackage{microtype,graphicx,booktabs,longtable,array,pdflscape}
\usepackage[font=small,labelfont=bf]{caption}
\usepackage[round,authoryear]{natbib}
\usepackage{xcolor,enumitem,setspace}
\usepackage{xr-hyper}
\usepackage[colorlinks=true,linkcolor=black,citecolor=black,urlcolor=blue]{hyperref}
\usepackage{fancyhdr}
\setlength{\headheight}{14pt}
\pagestyle{fancy}\fancyhf{}
\fancyhead[L]{\small lrcBART: historical information borrowing}
\fancyhead[R]{\small Manuscript draft}
\fancyfoot[C]{\thepage}
\renewcommand{\headrulewidth}{0.3pt}
\setstretch{1.12}
\setlength{\parindent}{1.2em}
\setlength{\parskip}{2pt}
\setlength{\emergencystretch}{2em}
\setlist[enumerate]{itemsep=3pt,topsep=4pt}
\newcommand{\E}{\operatorname{E}}
\newcommand{\Var}{\operatorname{Var}}
\newcommand{\Cov}{\operatorname{Cov}}
\newcommand{\ESS}{\operatorname{ESS}}
\newcommand{\N}{\mathrm{N}}
\newcommand{\IG}{\operatorname{IG}}
\newcommand{\ind}{\mathbf{1}}
%DIF 31c31
%DIF < \newcommand{\lrc}{lrcBART}
%DIF -------
\newcommand{\lrc}{LRC-BART} %DIF > 
%DIF -------
\newtheorem{theorem}{Theorem}
\newtheorem{lemma}{Lemma}[section]
\theoremstyle{remark}\newtheorem{remark}{Remark}[section]
\graphicspath{{tracked-source/figures/}}
\makeatletter
\def\input@path{{tracked-source/}}
\makeatother
\hypersetup{pdfauthor={},pdfsubject={Statistical methods manuscript draft}}
%DIF <  \fancyhead[R]{\small Main tables and figures}
%DIF -------
 \externaldocument[A-]{tracked-source/generated/appendix_refs}[appendix.pdf] %DIF > 
%DIF < \usepackage{tikz}
\fancyhead[R]{\small Main tables} %DIF > 
%DIF < \usetikzlibrary{arrows.meta}
%DIF PREAMBLE EXTENSION ADDED BY LATEXDIFF
%DIF CFONT PREAMBLE %DIF PREAMBLE
\RequirePackage{color}\definecolor{RED}{rgb}{1,0,0}\definecolor{BLUE}{rgb}{0,0,1} %DIF PREAMBLE
\DeclareOldFontCommand{\sf}{\normalfont\sffamily}{\mathsf} %DIF PREAMBLE
\providecommand{\DIFaddtex}[1]{{\protect\color{blue} \sf #1}} %DIF PREAMBLE
\providecommand{\DIFdeltex}[1]{{\protect\color{red} \scriptsize #1}} %DIF PREAMBLE
%DIF SAFE PREAMBLE %DIF PREAMBLE
\providecommand{\DIFaddbegin}{} %DIF PREAMBLE
\providecommand{\DIFaddend}{} %DIF PREAMBLE
\providecommand{\DIFdelbegin}{} %DIF PREAMBLE
\providecommand{\DIFdelend}{} %DIF PREAMBLE
\providecommand{\DIFmodbegin}{} %DIF PREAMBLE
\providecommand{\DIFmodend}{} %DIF PREAMBLE
%DIF FLOATSAFE PREAMBLE %DIF PREAMBLE
\providecommand{\DIFaddFL}[1]{\DIFadd{#1}} %DIF PREAMBLE
\providecommand{\DIFdelFL}[1]{\DIFdel{#1}} %DIF PREAMBLE
\providecommand{\DIFaddbeginFL}{} %DIF PREAMBLE
\providecommand{\DIFaddendFL}{} %DIF PREAMBLE
\providecommand{\DIFdelbeginFL}{} %DIF PREAMBLE
\providecommand{\DIFdelendFL}{} %DIF PREAMBLE
%DIF HYPERREF PREAMBLE %DIF PREAMBLE
\providecommand{\DIFadd}[1]{\texorpdfstring{\DIFaddtex{#1}}{#1}} %DIF PREAMBLE
\providecommand{\DIFdel}[1]{\texorpdfstring{\DIFdeltex{#1}}{}} %DIF PREAMBLE
%DIF COLORLISTINGS PREAMBLE %DIF PREAMBLE
\RequirePackage{listings} %DIF PREAMBLE
\RequirePackage{color} %DIF PREAMBLE
\lstdefinelanguage{DIFcode}{ %DIF PREAMBLE
%DIF DIFCODE_CFONT %DIF PREAMBLE
  moredelim=[il][\color{red}\scriptsize]{\%DIF\ <\ }, %DIF PREAMBLE
  moredelim=[il][\color{blue}\sffamily]{\%DIF\ >\ } %DIF PREAMBLE
} %DIF PREAMBLE
\lstdefinestyle{DIFverbatimstyle}{ %DIF PREAMBLE
	language=DIFcode, %DIF PREAMBLE
	basicstyle=\ttfamily, %DIF PREAMBLE
	columns=fullflexible, %DIF PREAMBLE
	keepspaces=true %DIF PREAMBLE
} %DIF PREAMBLE
\lstnewenvironment{DIFverbatim}{\lstset{style=DIFverbatimstyle}}{} %DIF PREAMBLE
\lstnewenvironment{DIFverbatim*}{\lstset{style=DIFverbatimstyle,showspaces=true}}{} %DIF PREAMBLE
%DIF END PREAMBLE EXTENSION ADDED BY LATEXDIFF

\hypersetup{hypertexnames=false}
\begin{document}
 \DIFaddbegin \begin{center}{\Large\bfseries Main tables}\DIFaddend \par\medskip
{ \DIFaddbegin \DIFadd{Joint treatment and source-discrepancy modeling with locally }\DIFaddend robust commensurate Bayesian additive regression trees } \end{center}
 \DIFaddbegin \begin{table}[!htbp]
\DIFaddendFL \centering
 \caption{ \DIFaddbeginFL \DIFaddFL{Descriptive joint treatment and control comparison across five Gaussian conditions}\DIFaddendFL . \DIFaddbeginFL \DIFaddFL{Each entry gives RMSE / 95\% interval coverage count / ATE convergence flags, with 40 datasets per condition. The flags count method--dataset groups with $\widehat R>1.05$ or bulk effective sample size below 400; they do not count individual chains. Comparative estimates remain provisional because of these diagnostics.}\DIFaddendFL }
 \DIFaddbeginFL \label{tab:joint_primary}
\setlength{\tabcolsep}{3.5pt}
\footnotesize
\begin{tabular}{lccc}
\DIFaddendFL \toprule
 \DIFaddbeginFL \DIFaddFL{Condition }\DIFaddendFL &  \DIFaddbeginFL \DIFaddFL{Joint LRC }\DIFaddendFL &  \DIFaddbeginFL \DIFaddFL{Joint source-BART }\DIFaddendFL &  \DIFaddbeginFL \DIFaddFL{Separated source-BART}\DIFaddendFL \\
\midrule
 \DIFaddbeginFL \DIFaddFL{Compatible }\DIFaddendFL &  \DIFaddbeginFL \DIFaddFL{0.1750 / 38 / 0 }\DIFaddendFL &  \DIFaddbeginFL \DIFaddFL{0.2133 / 39 / }\DIFaddendFL 1 &  \DIFaddbeginFL \DIFaddFL{0.2064 / 37 / }\DIFaddendFL 1 \\
 \DIFaddbeginFL \DIFaddFL{One-variable }\DIFaddendFL &  \DIFaddbeginFL \DIFaddFL{0.2068 / 38 / 3 }\DIFaddendFL &  \DIFaddbeginFL \DIFaddFL{0.2155 / 38 / }\DIFaddendFL 2  &  \DIFaddbeginFL \DIFaddFL{0.2105 / 37 / 0}\DIFaddendFL \\
 \DIFaddbeginFL \DIFaddFL{Two-variable }\DIFaddendFL &  \DIFaddbeginFL \DIFaddFL{0.2151 / 38 / 3 }\DIFaddendFL &  \DIFaddbeginFL \DIFaddFL{0.2147 / 38 / }\DIFaddendFL 2  &  \DIFaddbeginFL \DIFaddFL{0.2113 / 37 / 7}\DIFaddendFL \\
 \DIFaddbeginFL \DIFaddFL{Heterogeneous }\DIFaddendFL &  \DIFaddbeginFL \DIFaddFL{0.2092 / 38 / 8 }\DIFaddendFL &  \DIFaddbeginFL \DIFaddFL{0.2112 / 38 / 4 }\DIFaddendFL &  \DIFaddbeginFL \DIFaddFL{0.1917 / 39 / 5}\DIFaddendFL \\
 \DIFaddbeginFL \DIFaddFL{Null }\DIFaddendFL &  \DIFaddbeginFL \DIFaddFL{0.2298 / 36 / 4 }\DIFaddendFL &  \DIFaddbeginFL \DIFaddFL{0.2150 / 38 / 3 }\DIFaddendFL &  \DIFaddbeginFL \DIFaddFL{0.2100 / 38 / 8}\DIFaddendFL \\
 \DIFaddbeginFL \bottomrule
\end{tabular}
\end{table}
 \DIFaddend \clearpage
\begingroup
\setlength{\tabcolsep}{4pt}
\renewcommand{\arraystretch}{1.13}
\renewcommand{\arraystretch}{1.13}
\small
\begin{longtable}{lrr}
\caption{Characteristics of the harmonized multiple-myeloma analysis cohorts.}\label{tab:table5}\\
\toprule
Characteristic & EloKRd & UCMM \\ \midrule
\endfirsthead
\multicolumn{3}{l}{\tablename\ \thetable\ (continued)}\\
\toprule
Characteristic & EloKRd & UCMM \\ \midrule
\endhead
\midrule\multicolumn{3}{r}{Continued on next page}\\\endfoot
\bottomrule\endlastfoot
Participants & 30 & 200 \\
Age (years), median (IQR) & 62.8 (57.4, 70.2) & 60.6 (54.0, 66.9) \\
Male & 21 (70.0) & 109 (54.5) \\
Race &  &  \\
\quad White & 23 (76.7) & 131 (65.5) \\
\quad Black & 4 (13.3) & 58 (29.0) \\
\quad Other & 3 (10.0) & 11 (5.5) \\
Hispanic or Latino & 1 (3.3) & 10 (5.0) \\
High-risk cytogenetics & 15 (50.0) & 51 (25.5) \\
ASCT & 9 (30.0) & 174 (87.0) \\
Progression or death & 13 (43.3) & 127 (63.5) \\
All-cause death & 10 (33.3) & 43 (21.5) \\
\end{longtable}
\noindent\begin{minipage}{\linewidth}\footnotesize Entries are $n$ (\%) except where indicated. IQR, interquartile range; ASCT, autologous stem-cell transplantation. ASCT is a postinduction characteristic and is not presented as a pretreatment baseline variable. All analysis covariates were complete in both cohorts.\end{minipage}
\endgroup
 \clearpage
\begin{landscape}
\begingroup
\setlength{\tabcolsep}{4pt}
\renewcommand{\arraystretch}{1.13}
\renewcommand{\arraystretch}{1.04}
\fontsize{10}{11.5}\selectfont
\begin{longtable}{lrrrrr}
\caption{Five-year restricted mean survival time (RMST) in the application.}\label{tab:table6}\\
\toprule
Method & \shortstack{RMST ratio\\(95\% interval)} & \shortstack{RMST difference\\(95\% interval)} & \shortstack{EloKRd treatment\\RMST} & \shortstack{EloKRd hypothetical\\control RMST} & \shortstack{UCMM control\\RMST} \\ \midrule
\endfirsthead
\multicolumn{6}{l}{\tablename\ \thetable\ (continued)}\\
\toprule
Method & \shortstack{RMST ratio\\(95\% interval)} & \shortstack{RMST difference\\(95\% interval)} & \shortstack{EloKRd treatment\\RMST} & \shortstack{EloKRd hypothetical\\control RMST} & \shortstack{UCMM control\\RMST} \\ \midrule
\endhead
\midrule\multicolumn{6}{r}{Continued on next page}\\\endfoot
\bottomrule\endlastfoot
\multicolumn{6}{l}{\strut\textit{PFS}}\\*
KM & 1.040 (0.862, 1.254) & 0.143 (-0.548, 0.834) & 3.705 & --- & 3.562 \\
lrcBART & 0.997 (0.799, 1.229) & -0.009 (-0.768, 0.714) & 3.521 & 3.535 & --- \\
BART-CP & 0.988 (0.796, 1.215) & -0.042 (-0.760, 0.690) & 3.527 & 3.580 & --- \\
AFT-CP & 0.901 (0.708, 1.105) & -0.338 (-1.080, 0.335) & 3.069 & 3.408 & --- \\
hier. AFT & 1.062 (0.789, 1.467) & 0.192 (-0.834, 1.138) & 3.355 & 3.148 & --- \\
\multicolumn{6}{l}{\strut\textit{OS}}\\*
KM & 0.891 (0.776, 1.023) & -0.505 (-1.077, 0.068) & 4.122 & --- & 4.627 \\
lrcBART & 0.898 (0.765, 1.024) & -0.453 (-1.065, 0.101) & 3.995 & 4.456 & --- \\
BART-CP & 0.898 (0.769, 1.017) & -0.448 (-1.048, 0.073) & 3.985 & 4.439 & --- \\
AFT-CP & 0.809 (0.654, 0.951) & -0.814 (-1.521, -0.203) & 3.449 & 4.268 & --- \\
hier. AFT & 0.961 (0.761, 1.387) & -0.153 (-1.085, 1.073) & 3.807 & 3.953 & --- \\
\end{longtable}
\noindent\begin{minipage}{\linewidth}\footnotesize RMST and differences are in years. For KM, EloKRd treatment RMST and UCMM control RMST are observed cohort estimates. For adjusted methods, EloKRd treatment and hypothetical control RMSTs are model-based estimates standardized to EloKRd covariates; UCMM control RMST is not reported. A dash denotes an inapplicable estimate. KM is an unadjusted comparison with confidence intervals; model-based intervals are equal-tailed credible intervals. Arm summaries are posterior medians, except KM estimates. The median of a ratio or difference need not equal the ratio or difference of marginal medians. Primary lrcBART uses $H_f=10$, $H_g=5$, $w=1$, requested ESS 30. The hier. AFT comparator uses hierarchical discrepancy-variance prior scale $s_\tau^2=0.05$, where $\tau_\alpha^2,\tau_\beta^2\sim\operatorname{IG}(3/2,3s_\tau^2/2)$.\end{minipage}
\endgroup
\end{landscape}
 \clearpage
\begingroup
\setlength{\tabcolsep}{4pt}
\renewcommand{\arraystretch}{1.0}
\fontsize{9.5}{11}\selectfont
\begin{longtable}{lrlllllrr}
\caption{UCMM-informed prior ESS for the 30 EloKRd patient covariate profiles, with UCMM covariate-combination counts (PFS).}\label{tab:table7}\\
\toprule
Age group & Age & Sex & Race & \shortstack{Hispanic/\\Latino} & \shortstack{Cytogenetic\\risk} & ASCT & \shortstack{UCMM\\cohort $n$} & Prior ESS \\ \midrule
\endfirsthead
\multicolumn{9}{l}{Table \thetable\ (continued)}\\
\toprule
Age group & Age & Sex & Race & \shortstack{Hispanic/\\Latino} & \shortstack{Cytogenetic\\risk} & ASCT & \shortstack{UCMM\\cohort $n$} & Prior ESS \\ \midrule
\endhead
\midrule\multicolumn{9}{r}{Continued on next page}\\\endfoot
\bottomrule\endlastfoot
$<50$ & 42.8 & Male & White & No & High & Yes & 0 & 11.0 \\
$<50$ & 45.4 & Male & White & No & High & No & 0 & 6.0 \\
$<50$ & 48.4 & Male & Black & No & Standard & No & 0 & 6.7 \\
\addlinespace[3pt]
50--59 & 50.7 & Female & White & No & High & No & 0 & 4.7 \\
50--59 & 51.7 & Male & Other & No & Standard & Yes & 1 & 7.1 \\
50--59 & 55.7 & Male & White & No & Standard & No & 2 & 8.3 \\
50--59 & 56.6 & Male & White & No & Standard & Yes & 22 & 21.2 \\
50--59 & 57.2 & Male & White & No & Standard & Yes & 22 & 21.2 \\
50--59 & 58.1 & Female & White & No & Standard & Yes & 17 & 15.8 \\
50--59 & 59.0 & Female & Black & No & High & No & 0 & 5.1 \\
50--59 & 59.4 & Male & White & No & Standard & Yes & 22 & 21.7 \\
50--59 & 59.5 & Male & White & No & High & Yes & 6 & 13.1 \\
50--59 & 59.7 & Female & White & No & Standard & No & 0 & 7.0 \\
\addlinespace[3pt]
60--69 & 61.6 & Female & White & No & High & No & 0 & 5.5 \\
60--69 & 62.4 & Male & White & No & High & No & 2 & 7.2 \\
60--69 & 63.2 & Male & Other & No & High & No & 0 & 4.3 \\
60--69 & 65.6 & Male & White & No & Standard & Yes & 18 & 24.5 \\
60--69 & 67.4 & Male & White & No & High & No & 2 & 7.2 \\
60--69 & 67.7 & Female & White & No & Standard & No & 0 & 7.6 \\
60--69 & 67.9 & Male & Other & Yes & Standard & No & 0 & 3.8 \\
60--69 & 69.7 & Male & White & No & High & No & 2 & 7.2 \\
\addlinespace[3pt]
$\geq70$ & 70.0 & Male & White & No & Standard & No & 4 & 9.4 \\
$\geq70$ & 70.3 & Male & White & No & High & No & 0 & 7.1 \\
$\geq70$ & 70.4 & Male & White & No & High & Yes & 2 & 13.5 \\
$\geq70$ & 70.4 & Male & White & No & High & No & 0 & 7.1 \\
$\geq70$ & 70.9 & Male & White & No & High & No & 0 & 7.1 \\
$\geq70$ & 72.1 & Male & White & No & Standard & No & 4 & 9.5 \\
$\geq70$ & 76.1 & Female & White & No & High & No & 1 & 5.6 \\
$\geq70$ & 77.0 & Female & Black & No & Standard & No & 3 & 7.2 \\
$\geq70$ & 81.2 & Female & Black & No & Standard & No & 3 & 7.0 \\
\end{longtable}
\noindent\begin{minipage}{\linewidth}\footnotesize Each row is one EloKRd patient covariate profile. Age groups are $<50$, 50--59, 60--69 and $\geq70$ years. Age is rounded to one decimal year for display; calculations use exact age. UCMM cohort $n$ is the number of patients among all 200 triplet-treated UCMM controls matching that row's age group, sex, race, Hispanic/Latino status, cytogenetic risk and ASCT. Counts repeat when EloKRd rows share the same displayed combination and should not be summed down the table. Prior ESS is informed by the full 200-patient triplet UCMM cohort and is expressed in uncensored normal-reference units for the conditional mean log-PFS at the same profile. The same residual variance and globally selected PFS scale are used for all profiles. Values do not measure outcome compatibility and do not sum or average to overall ESS. ASCT is postinduction.\end{minipage}
\endgroup
  \end{document}
